% ============================================================
% Appendix D: Further Reading
% ============================================================

\chapter*{Suggestions for Further Reading}
\addcontentsline{toc}{chapter}{Suggestions for Further Reading}

\section*{General References}

\begin{itemize}
    \item Constance Reid, \textit{Hilbert}, Springer-Verlag, 1970. A classic biography of David Hilbert, readable and informative.

    \item Jeremy Gray, \textit{The Hilbert Challenge}, Oxford University Press, 2000. A detailed history of Hilbert's problems and their influence on 20th-century mathematics.

    \item Barry Mazur, ``The Hilbert Problems,'' in \textit{Mathematics: Frontiers and Perspectives}, AMS, 2000. An essay on the philosophical implications of Hilbert's problems.
\end{itemize}

\section*{Problem-Specific References}

\begin{description}
    \item[Problem 1] Paul Cohen, \textit{Set Theory and the Continuum Hypothesis}, Benjamin, 1966. The original monograph on forcing and independence.

    \item[Problem 2] Raymond Smullyan, \textit{G\"odel's Incompleteness Theorems}, Oxford University Press, 1992. An accessible introduction. See also George Boolos and Richard Jeffrey, \textit{Computability and Logic}, Cambridge University Press.

    \item[Problem 8] Tom\"a\v{s} Oliveira e Silva, Siegfried Herzog, and Christian Pelayo, ``Computational verification of the Ternary Goldbach Conjecture up to $8 \times 10^{13}$,'' \textit{Mathematics of Computation}, 2014. For computational approaches.

    \item[Problem 10] Martin Davis, ``Hilbert's Tenth Problem Is Unsolvable,'' \textit{The American Mathematical Monthly}, 1973. A clear expository account.

    \item[Problem 13] Ian Stewart, \textit{Galois Theory}, Chapman \& Hall, 2015. The definitive undergraduate textbook on Galois theory.

    \item[Problem 16] Claire Voisin, \textit{Hodge Theory and Complex Algebraic Geometry}, Cambridge University Press, 2002--2003. An advanced but rewarding treatment.

    \item[Problem 22] L. Ahlfors, \textit{Complex Analysis}, McGraw-Hill, 1979. Chapter 8 covers the uniformization theorem.
\end{description}

\section*{Textbooks for Further Study}

\begin{itemize}
    \item Michael Artin, \textit{Algebra}, Pearson, 2011. A standard undergraduate algebra text covering group theory, rings, and Galois theory.

    \item John M. Lee, \textit{Introduction to Smooth Manifolds}, Springer, 2013. For the topology and geometry underlying several of Hilbert's problems.

    \item Lawrence C. Evans, \textit{Partial Differential Equations}, AMS, 2010. For the PDE theory behind Problems 19 and 20.

    \item Serge Lang, \textit{Algebraic Number Theory}, Springer, 2000. For number theory beyond the undergraduate level.

    \item Miles Reid, \textit{Undergraduate Algebraic Geometry}, Cambridge University Press, 1988. An accessible introduction to the geometry underlying Problems 14--16.
\end{itemize}

\section*{Online Resources}

\begin{itemize}
    \item The Clay Mathematics Institute's description of the Millennium Problems: \url{https://www.claymath.org/millennium-problems/}

    \item The Wikipedia article on Hilbert's problems provides a good overview with references.

    \item \url{https://mathshistory.st-andrews.ac.uk/}: The MacTutor History of Mathematics archive, with biographical and historical information on many mathematicians mentioned in this book.
\end{itemize}
